\documentclass[10pt,a4paper]{amsart}
\usepackage[margin=19mm]{geometry}
\usepackage{amsmath,amssymb,mathtools}
\usepackage[scr=rsfs,cal=euler]{mathalfa}
\usepackage{xfrac}
\usepackage[unicode,hidelinks,bookmarks=false]{hyperref}
\newcommand{\ie}{i.e.\ }
\newcommand{\cf}{cf.\ }
\newcommand{\resp}{respectively\ }
\newcommand{\Subsection}{\subsection}
\providecommand{\nobreakdash}{\nobreak\hbox{-}}
\setlength{\emergencystretch}{2em}
\multlinegap0pt
\usepackage{bm}
\usepackage{bbm}
\usepackage{dsfont}
\usepackage{xspace}
\usepackage{cancel}
\usepackage{mathtools}
\usepackage{amsthm}
\makeatletter
\@ifundefined{theorem}{\newtheorem{theorem}{Theorem}}{}
\@ifundefined{lemma}{\newtheorem{lemma}[theorem]{Lemma}}{}
\makeatother
\makeatletter
\newcommand{\Arg}{\operatorname{Arg}}
\renewcommand{\Re}{\operatorname{Re}}
\expandafter\ifx\csname pf\endcsname\relax
\newenvironment{pf}
   {\par\noindent{\it Proof of}\hskip 0.5em\ignorespaces}
\fi
\makeatother
\title[An extension of Ramanujan-Guinand identity for the Dedekind ]{An extension of Ramanujan-Guinand identity for the Dedekind zeta function and a new formula for $\zeta\_\mathbb{F}(2)$ and $\zeta\_\mathbb{F}(2m+1)$}
\author{Diksha Rani Bansal, Bibekananda Maji}
\date{Source: arXiv:2608.07094v1 (CC BY 4.0)}
\begin{document}
\maketitle

\noindent\emph{Source and licence.} An extension of Ramanujan-Guinand identity for the Dedekind zeta function and a new formula for $\zeta_\mathbb{F}(2)$ and $\zeta_\mathbb{F}(2m+1)$. Version arXiv:2608.07094v1 (CC BY 4.0), distributed under the Creative Commons Attribution 4.0 International licence, as recorded in the arXiv licence metadata for this exact version and shown on the abstract page https://arxiv.org/abs/2608.07094v1. Full rights notice: licenses/arxiv-2608.07094-CC-BY-4.0.txt.\par\smallskip

\noindent\textit{Editorial scope.} Counted unit: the Lemma whose source label is \texttt{Lemma for Mellin trans} in the article (source0.tex lines 483--489, source label \texttt{Lemma for Mellin trans}), delivered together with the article's own complete original proof of it (source0.tex lines 490--513, one \texttt{proof} environment of 24 lines). No mathematics is written, repaired or re-derived: statement and proof are the article's own text line for line. Required context retained uncounted: Steen fn (lines 257--259); Short notation (lines 271--273); Legendre's Duplication Formula (lines 440--442). Every definition, display and label of those lines is retained unchanged. Omitted from this package and not redistributed: 1--256, 260--270, 274--439, 443--482, 514--990. Nothing inside a retained excerpt is omitted or altered: every retained body is delivered exactly as the article's lines are written, so no omission receipt of any kind accompanies this package. Citations: the retained excerpts carry no citation command at all, so no bibliography entry of the article is transcribed and no reference is redistributed. Reference adaptation: none was needed and none was made; every reference inside the delivered unit resolves inside it, so no unresolved reference or citation is produced. Printed slips kept literal: no printed word, symbol, hypothesis or number of the counted statement or of its proof was altered. Layout and class: the article's own class is \texttt{amsart} with packages including amsmath, amssymb, amsthm, amsfonts, url, mathtools, hyperref, tikz, array; the wrapper is the lane's pinned \texttt{amsart} editorial wrapper at 10pt on A4, which restates the theorem-like environments the retained text needs and adds the article's own macro definitions that the retained text needs (\texttt{\textbackslash Arg}, \texttt{\textbackslash Re}), with extra wrapper packages (bm, bbm, dsfont, xspace, cancel, mathtools). The delivered numbering is the unit's own. Assets: no retained span contains a figure environment, a graphics-inclusion command, a PDF-page inclusion, a table or a TikZ picture; no image file is copied into this package, the package declares no asset dependency and no image file is redistributed with it. Licence: version arXiv:2608.07094v1 of this article is distributed under the Creative Commons Attribution 4.0 International licence, as recorded in the arXiv licence metadata for this exact version and shown on the abstract page https://arxiv.org/abs/2608.07094v1. A full rights notice is delivered at licenses/arxiv-2608.07094-CC-BY-4.0.txt. Characters: every retained excerpt is pure ASCII, so the delivered engine, XeLaTeX with its default Latin Modern text font, reports no missing character.
\begin{align}\label{Steen fn}
V(x| a_1, a_2, \ldots, a_n) := \frac{1}{2\pi i}\int\limits_{(c)}\prod_{j=1}^{n}\Gamma(s + a_j)x^{-s}\text{d}s, \quad |\Arg(x)| < \frac{\pi n}{2}, 
\end{align}

\begin{align}\label{Short notation}
V_{\mathbb{F},z}(x) := V\left(x\bigg|\left(0\right)_{r_1 + r_2}, \left(\frac{z}{2}\right)_{r_1 + r_2}, \left(\frac{1}{2}\right)_{r_2}, \left(\frac{1+z}{2}\right)_{r_2}\right),
\end{align}

\begin{align}\label{Legendre's Duplication Formula}
\Gamma(s)\Gamma\left(s + \frac{1}{2}\right)  &= 2^{1-2s}\sqrt{\pi}\Gamma(2s),  \quad \Gamma(s+1) = s\Gamma(s).
\end{align}

\begin{lemma}\label{Lemma for Mellin trans}
The following identity holds for $\Re(s) > \max\{0, -\Re(z)\}$,
\begin{align*}
\left(\frac{D}{4^{r_2}\pi^{d}}\right)^{s}\Gamma^{r_1}\left(\frac{s}{2}\right)\Gamma^{r_1}\left(\frac{s+z}{2}\right)\Gamma^{r_2}(s)\Gamma^{r_2}(s+z)
&= \frac{2^{r_2 z + 1}}{(4\pi)^{r_2}}\int\limits_{0}^\infty x^{s-1}V_{\mathbb{F},z}\left(\left(\frac{\pi^{d} x}{D}\right)^2\right) \mathrm{d}x.
\end{align*}
\end{lemma}

\begin{proof}
From \eqref{Short notation} and the definition \eqref{Steen fn} of the Steen function, we have
\begin{align*}
V_{\mathbb{F},z}\left(x\right)
&= \frac{1}{2\pi i}\int\limits_{(c)}\left[\Gamma(s)\Gamma\left(s+\frac{z}{2}\right)\right]^{r_1+r_2}\!\left[\Gamma\left(s+\frac{1}{2}\right)\Gamma\left(s+\frac{1+z}{2}\right)\right]^{r_2}\!x^{-s} \textrm{d}s,
\end{align*}
where $c > \max\left\{0, -\frac{\Re(z)}{2}\right\}$.
We apply the duplication formula \eqref{Legendre's Duplication Formula} for the gamma function in the last equation and subsequently perform a change of variable $s \mapsto \frac{s}{2}$ to arrive at
\begin{align*}
V_{\mathbb{F},z}\left(x\right)
&= \frac{(4\pi)^{r_2}}{2^{r_2 z} 2\pi i}{\displaystyle\int\limits_{(c)}\!\left[\Gamma(s)\Gamma\left(s+\frac{z}{2}\right)\right]^{r_1}\!\left[\Gamma(2s)\Gamma(2s+z)\right]^{r_2}(4^{2r_2}x)^{-s}\textrm{d}s}\\
&= \frac{(4\pi)^{r_2}}{2^{r_2 z + 1} 2\pi i}{\displaystyle\int\limits_{(2c)}\!\Gamma^{r_1}\left(\frac{s}{2}\right)\Gamma^{r_1}\left(\frac{s+z}{2}\right)\Gamma^{r_2}(s)\Gamma^{r_2}(s+z)(4^{r_2}\sqrt{x})^{-s}\textrm{d}s}.
\end{align*}
We now substitute $x = \left(\frac{\pi^{d} x}{D}\right)^2$ to have
\begin{align}\label{Inverse Mellin}
V_{\mathbb{F},z}\left(\left(\frac{\pi^{d} x}{D}\right)^2\right)
&= \frac{(4\pi)^{r_2}}{2^{r_2 z + 1}2\pi i}{\displaystyle\int\limits_{(2c)}\!\left(\frac{D}{4^{r_2}\pi^{d}}\right)^{s}\Gamma^{r_1}\left(\frac{s}{2}\right)\Gamma^{r_1}\left(\frac{s+z}{2}\right)\Gamma^{r_2}(s)\Gamma^{r_2}(s+z)x^{-s}\textrm{d}s}.
\end{align}
Hence, the Mellin transform for the product of the gamma factors in the integrand above is given by
\begin{align*}
\left(\frac{D}{4^{r_2}\pi^{d}}\right)^{s}\Gamma^{r_1}\left(\frac{s}{2}\right)\Gamma^{r_1}\left(\frac{s+z}{2}\right)\Gamma^{r_2}(s)\Gamma^{r_2}(s+z)
&= \frac{2^{r_2 z + 1}}{(4\pi)^{r_2}}\int\limits_{0}^\infty x^{s-1}V_{\mathbb{F},z}\left(\left(\frac{\pi^{d} x}{D}\right)^2\right) \mathrm{d}x.
\end{align*}
\end{proof}

\end{document}
